\begin{lstlisting}[language=diff,basicstyle=\ttfamily\scriptsize,breaklines=true,caption={tsvc/s121 — FASTER (e1_bare_a)}]
--- original.c
+++ final.c
@@ -134,13 +134,38 @@
 static real_t kernel_s121(void)
 {
     int j;
+    /* Allocate temporary buffer to hold old values of a.
+     * This breaks loop-carried dependence by letting all iterations
+     * read from a constant copy instead of writing-then-reading in sequence. */
+    real_t *a_old = (real_t*)malloc((size_t)LEN_1D * sizeof(real_t));
+    if (!a_old) {
+        fprintf(stderr, "out of memory in kernel_s121\n");
+        return (real_t)0;
+    }
+
     for (int nl = 0; nl < R; nl++) {
+        /* Copy current a to a_old for this iteration.
+         * j is the loop variable (automatically private).
+         * Loop is parallelizable since each thread writes to a different element. */
+        #pragma omp parallel for
+        for (long i = 0; i < LEN_1D; i++) {
+            a_old[i] = a[i];
+        }
+
+        /* Main computation loop: read from a_old (constant), write to a (independent).
+         * j: private variable, each thread computes its own value (j = i+1).
+         * Loop is now parallelizable since iterations are independent:
+         * all read from the same a_old copy, write to different a[i] locations. */
+        #pragma omp parallel for private(j)
         for (int i = 0; i < LEN_1D-1; i++) {
             j = i + 1;
-            a[i] = a[j] + b[i];
+            a[i] = a_old[j] + b[i];
         }
+
         pb_mix(nl);
     }
+
+    free(a_old);
     return (real_t)0;
 }
 
\end{lstlisting}
